\documentclass[10pt,a4paper]{amsart}
\usepackage[margin=19mm]{geometry}
\usepackage{amsmath,amssymb,mathtools}
\usepackage[scr=rsfs,cal=euler]{mathalfa}
\usepackage{xfrac}
\usepackage[unicode,hidelinks,bookmarks=false]{hyperref}
\newcommand{\ie}{i.e.\ }
\newcommand{\cf}{cf.\ }
\newcommand{\resp}{respectively\ }
\newcommand{\Subsection}{\subsection}
\providecommand{\nobreakdash}{\nobreak\hbox{-}}
\setlength{\emergencystretch}{2em}
\multlinegap0pt
\usepackage{dsfont}
\usepackage{xcolor}
\usepackage{csquotes}
\usepackage{mathtools}
\usepackage{amssymb}
\usepackage{bm}
\newtheorem{lemma}{Lemma}
\DeclareMathOperator{\Id}{Id}
\DeclareMathOperator{\trace}{trace}
\title{Required Number of Points in $L_2$~Marcinkiewicz--Zygmund Inequalities}
\author{Felix Bartel}
\date{Source: arXiv:2608.25886v1 (CC BY 4.0)}
\begin{document}
\maketitle

\noindent\emph{Source and licence.} Required Number of Points in $L_2$~Marcinkiewicz--Zygmund Inequalities. Version arXiv:2608.25886v1 (CC BY 4.0), distributed by arXiv under the Creative Commons Attribution 4.0 International licence, http://creativecommons.org/licenses/by/4.0/, as recorded in the arXiv licence metadata for this exact version and shown on the abstract page https://arxiv.org/abs/2608.25886v1. Full licence notice: \texttt{licenses/arxiv-2608.25886-CC-BY-4.0.txt}.\par\smallskip

\noindent\textit{Editorial scope.} Counted unit: the article's lemma edgeframespectrum (source0.tex lines 1782--1798). The article's complete original proof of it is delivered with it (source0.tex lines 1800--1869, one \texttt{proof} environment of 70 lines). No mathematics is written, repaired or re-derived: the statement and the proof are the article's own lines, in the article's own notation, and no retained line is substituted or re-flowed.  Retained uncounted as the source-local context the proof needs: lines 1366--1368; lines 1774--1777.  Nothing was omitted from any retained span, so the package carries no omission record.  No retained line contains a citation, so the wrapper carries no bibliography and no bibliography entry.  No retained line contains a figure environment, an image-inclusion command or drawing code, so no image is delivered and the package declares no asset dependency.  Editorial formatting only, in addition: an AMS \texttt{amsart} preamble that restates the article's own declarations and macros used by the retained lines, namely the environments \texttt{lemma} on separate counters, together with the standard packages those lines need.  The retained environments print in this document as Lemma 1 in order of appearance, whereas the article prints the same items with the numbering of its own full document.  The complete native source of the article as distributed in the arXiv e-print for this exact version is bundled unchanged as \texttt{sources/208601/source0.tex}.
\begin{equation}\label{eq:gram}
    \bm G_\Phi = \Big[\trace(Q_iQ_j)\Big]_{i,j=1}^{N} \in\mathds R^{N\times N} .
\end{equation}

\begin{equation}\label{eq:edgeframe}
    \bm\varphi_e
    =\frac{\bm e_u-\bm e_v}{\sqrt2}\in \mathcal H_{\mathcal K}.
\end{equation}

\begin{lemma}[Spectrum of the centered edge projectors]\label{edgeframespectrum}
    The vectors $\{\bm\varphi_e : e\in\binom{V}{2}\}$ from \eqref{eq:edgeframe} form a UNTF for $\mathcal H_{\mathcal K}$.
    The centered-projector Gram matrix $\bm G_{\mathcal K}$ from \eqref{eq:gram} has eigenvalue zero on $\bm 1_N$, and, on $\bm 1_N^\perp$, the eigenvalues
    \begin{equation*}
        \frac {|V|}4 \text{ with multiplicity }|V|-1
        \quad\text{and}\quad
        \frac12 \text{ with multiplicity }\frac{|V|(|V|-3)}2.
    \end{equation*}
    Consequently,
    \begin{equation*}
        \gamma_{\mathcal K}
        =\begin{cases}
            3/4,&|V|=3,\\
            1/2,&|V|\ge4.
        \end{cases}
    \end{equation*}
\end{lemma}

\begin{proof}
    For the projector $P_e\colon\mathcal H_{\mathcal K}\to\mathcal H_{\mathcal K}, \bm a\mapsto \langle\bm a,\bm\varphi_e\rangle\bm\varphi_e$, we obtain
    \begin{equation*}
        \sum_{e\in\binom V2}P_e
        =\frac12\bm L_{\mathcal K}\big|_{\mathcal H_{\mathcal K}}
        =\frac {|V|}2\Id_{\mathcal H_{\mathcal K}}
        =\frac Nm\Id_{\mathcal H_{\mathcal K}}.
    \end{equation*}
    Since the edge vectors $\bm\varphi_e$ are unit-norm by definition, they form a UNTF.

    Let $\bm A_{\mathrm L} \in\mathds R^{N\times N}$ be the adjacency matrix of the line graph of $\mathcal K$.
    Let
    \begin{equation*}
        \bm B\in\mathds R^{|V|\times N}
        \quad\text{with entries}\quad
        b_{v,e} = \begin{cases}1 & \text{for } v\in e\,,\\0& \text{otherwise.}\end{cases}
    \end{equation*}
    be the unsigned vertex-edge incidence matrix of $\mathcal K$.
    Because every edge contains two vertices and two distinct edges contribute one precisely when they meet, we have the relation
    \begin{equation*}
        \bm B^\top\bm B
        = \Big[ \sum_{v\in V} \delta_{v\in e}\delta_{v\in f} \Big]_{e,f\in \binom{V}{2}}
        = 2\bm I_N+\bm A_{\mathrm L} .
    \end{equation*}
    On the other hand, we have
    \begin{equation*}
        \bm B\bm B^\top
        = \Big[ \sum_{e\in\binom{V}{2}} \delta_{v\in e}\delta_{u\in e} \Big]_{u,v\in V}
        = (|V|-2)\bm I_{|V|}+\bm 1_{|V|}\bm 1_{|V|}^\top .
    \end{equation*}
    The latter matrix has eigenvalues $2(|V|-1)$ with multiplicity one and $|V|-2$ with multiplicity $|V|-1$.
    Since $\bm B^\top\bm B$ has the same nonzero eigenvalues, the spectrum of the line-graph matrix $\bm A_{\mathrm L}$ is
    \begin{center}
    \begin{tabular}{c|ccc}
        eigenvalue & $2(|V|-2)$ & $|V|-4$ & $-2$ \\
        \hline
        multiplicity & $1$ & $|V|-1$ & $N-|V|$
    \end{tabular}
    \end{center}

    Now we relate the spectrum of $\bm A_{\mathrm L}$ to the centered-projector Gram matrix $\bm G_{\mathcal K}$.
    For the Gram matrix of the projectors we have the entries
    \begin{equation*}
        \trace(P_eP_f)
        = |\langle\bm\varphi_e,\bm\varphi_f\rangle|^2
        = \begin{cases} 1 & \text{for }e=f \,, \\ 1/4 & \text{if two edges } e\neq f \text{ meet in one vertex,}\\ 0 & \text{if the edges } e\text{ and } f \text{ are disjoint.}\end{cases}
    \end{equation*}
    Therefore $[\trace(P_eP_f)]_{e,f\in\binom{V}{2}} = \bm I_{N} + \frac 14 \bm A_{\mathrm L}$ has the eigenvalues
    \begin{center}
    \begin{tabular}{c|ccc}
        eigenvalue & $|V|/2$ & $|V|/4$ & $1/2$ \\
        \hline
        multiplicity & $1$ & $|V|-1$ & $N-|V|$
    \end{tabular}
    \end{center}
    Because every edge meets $2(|V|-2)$ other edges, we obtain
    \begin{equation*}
        \bm A_{\mathrm L}\bm 1_{N} = 2(|V|-2)\bm 1_{N}
        \quad\Leftrightarrow\quad
        [\trace(P_eP_f)]_{e,f\in\binom{V}{2}}\bm 1_{N}
        = \frac{|V|}{2}\bm 1_{N} ,
    \end{equation*}
    i.e., this determines the one-dimensional eigenspace for the eigenvalue $|V|/2$.
    To obtain the Gram matrix for the centered projectors, we have to subtract $m^{-1}\bm 1_N\bm 1_N^\top$
    \begin{equation*}
        \bm G_{\mathcal K}
        = \bm I_{N} + \frac 14 \bm A_{\mathrm L} - \frac 1m\bm 1_{N}\bm 1_{N}^\top .
    \end{equation*}
    This cancels the eigenvalue ${|V|}/2$ in the direction $\bm 1_N$ and leaves the other eigenvalues unchanged, proving the assertion.
\end{proof}

\end{document}
